% Fragmento público generado desde propietarios íntegros.
% Residencia: 52.11.2.
% Fuente: ../incoming/radion_monografia_20260827/manuscrito/sections/03_cierre_exacto.tex:286-309
\subsubsection{Caractérisation structurelle de \(\pi\) par compatibilité des sections}

Réordonner \eqref{eq:cierre-unitario} donne
\begin{equation}
\boxed{
2\pih=
\frac{1+\Phi_T-\epsone}
{\frac5{36}+\frac{25}{9}\alphah}.}
\label{eq:pi-radion}
\end{equation}
Cette égalité peut être appelée une nouvelle définition de \(\pi\) au sein de la
carte du radion si l'on entend \emph{définition structurelle} comme une
caractérisation exacte par compatibilité des sections. Ce n'est pas un second
générateur indépendant : \(\Delta_4\) et l'opérateur de retenue reçoivent déjà la
section coinductive de \(\pi\). La nouveauté réside dans le fait que le tour est
caractérisé par le raccord entre publication directe, torsion et mémoire.

\begin{narrativebox}
L'équation \eqref{eq:pi-radion} ne dit pas que nous ignorons d'abord \(\pi\) et
le devinons en fermant un angle. Elle dit que le \(\pi\) généré par la connexion
nonadique est exactement celui qui rend compatibles deux lectures internes du
même état. C'est une identité de reconnaissance intrinsèque, non un
étalonnage métrologique.
\end{narrativebox}
